% =====================================================================
\section{Results on the X3}
\label{sec:x3res}
% =====================================================================

All experiments run on the X3 surrogate (\cref{sec:x3}) at 60\,Hz with one sample of delay, with a tiny asymmetric
torque of $2\times10^{-4}$\,N\,m always applied so that a perfectly symmetric simulation cannot hide an instability.
Numbers: \code{docs/results/x3\_metrics.md}; script: \code{x3\_sim/run\_x3\_experiments.py}.

\begin{center}\small
\begin{tabularx}{\linewidth}{@{}l X@{}}
\toprule
\textbf{Exp.} & \textbf{Question}\\
\midrule
E0/E1 & Does the port of v1 crash like Gazebo did, and which fault causes it?\\
E2 & Does the corrected paper controller hover for 180\,s inside its envelopes?\\
E3 & With the paper's disturbance (scaled to what the X3 can produce), constrained vs.\ conventional PD?\\
E3b & What happens at the full paper disturbance?\\
E4 & Can it track a trajectory (2\,m\,$\times$\,2\,m square at 2\,m)?\\
E5 & How sensitive is it to control rate, delay and time-stamp jitter?\\
\bottomrule
\end{tabularx}
\end{center}

% ---------------------------------------------------------------------
\subsection{E0/E1: the crash of v1, reproduced and explained}
% ---------------------------------------------------------------------

\resfig{x3_E1_fault_isolation.png}{\linewidth}{Fault isolation: v1 with each fix switched on separately (60\,s, $\times$ = crash).}{fig:e1}

\begin{table}[H]
\centering\small
\caption{E1 fault isolation (60\,Hz).}
\label{tab:e1}
\begin{tabular}{@{}lcccc@{}}
\toprule
\textbf{Variant} & \textbf{crashed} & \textbf{crash time} & \textbf{max altitude} & \textbf{RMS tilt after 20\,s}\\
\midrule
v1 as is & yes & 5.07\,s & 0.24\,m & --\\
v1 + z-rate fix & yes & 5.68\,s & 1.06\,m & --\\
v1 + mixer fix & yes & 5.07\,s & 0.24\,m & --\\
v1 + mixer + z-rate fixes & no & -- & 2.00\,m & $2.03^\circ$\\
v1 + mixer + z-rate + body-rate fixes & no & -- & 1.97\,m & $0.007^\circ$\\
v1 as is at 250\,Hz & yes & 14.67\,s & 0.15\,m & --\\
\bottomrule
\end{tabular}
\end{table}

\resfig{x3_E0_v1_failure.png}{\linewidth}{v1 with only the z-rate fixed: the behaviour we saw in Gazebo. It climbs, then the
attitude diverges because pitch has no authority, the commanded roll torque saturates, and the vehicle crashes at 5.7\,s.}{fig:e0}

\begin{projbox}[conclusion of E1]
Two faults are each sufficient to crash the vehicle: the altitude-rate bug F14 (it never climbs above 0.24\,m) and the
mixer F9 (with F14 fixed it climbs to 1.06\,m and then tips over, exactly as observed in Gazebo). Fixing both gives a
stable flight; the remaining slow attitude limit cycle ($2^\circ$ RMS) comes from the lagging body-rate estimate F11 and
disappears with the exact estimate. A faster loop (250\,Hz) does not save v1: the faults are structural, not a matter of
rate.
\end{projbox}

% ---------------------------------------------------------------------
\subsection{E2: 180\,s hover with the paper's law}
% ---------------------------------------------------------------------

\resfig{x3_E2_hover_180s.png}{\linewidth}{E2, corrected controller (paper barrier form), 180\,s: altitude, horizontal position,
attitude error and body rate inside their envelopes, rotor commands and $\that$.}{fig:e2}

\begin{table}[H]
\centering\small
\caption{E2 metrics (after the 8\,s take-off ramp where stated).}
\label{tab:e2}
\begin{tabular}{@{}lc@{\hspace{1.5cm}}lc@{}}
\toprule
\textbf{Metric} & \textbf{Value} & \textbf{Metric} & \textbf{Value}\\
\midrule
crashed & no & time outside $\rho_q$ / $\rho_\omega$ & 0\,\% / 0\,\%\\
altitude RMS error ($t\ge10$\,s) & 0.02\,mm & $\max|\varepsilon_1|$ / $\max|\varepsilon_2|$ & 0.0023 / 0.0048\\
altitude max error ($t\ge10$\,s) & 0.24\,mm & mixer saturated & 0\,\% of the time\\
horizontal drift (max) & 0.21\,mm & rotor speeds ($t\ge10$\,s) & 659.9--660.7\,rad/s\\
tilt RMS ($t\ge10$\,s) & $0.003^\circ$ & $\that$ at 180\,s & 0.0012\,N\,m\\
\bottomrule
\end{tabular}
\end{table}

\begin{warn}[How to read these very small numbers]
The surrogate has no sensor noise and only the tiny test torque, so these are the errors of the control loop itself
(it is stable, well damped and has no limit cycle), not predictions of the errors in Gazebo or on hardware. The rotor
speeds sit at the hover value $660.4$\,rad/s computed by hand in the Mathematical Models (block 10).
\end{warn}

% ---------------------------------------------------------------------
\subsection{E3: disturbance rejection, constrained vs.\ conventional PD}
% ---------------------------------------------------------------------

The paper's disturbance has a yaw amplitude of 0.15\,N\,m, while the X3 can produce at most 0.106\,N\,m of yaw torque at
hover. E3 therefore uses $0.5\times$ the paper's $\vd(t)$, applied from $t=15$\,s to 180\,s (E3b uses $1\times$).

\resfig{x3_E3_disturbance_rates.png}{\linewidth}{E3, body rates against the rate envelope ($\rho_{\omega T}=0.05$\,rad/s), first
60\,s. Black: outside. The paper form touches the \emph{lower} yaw wall, the symmetric form never leaves, conventional PD
leaves at both walls on every yaw oscillation.}{fig:e3}

\resfig{x3_E3_violation_bars.png}{0.8\linewidth}{E3, percentage of time outside each envelope (15--180\,s).}{fig:e3bars}

\begin{table}[H]
\centering\small
\caption{E3 metrics, disturbance $0.5\times$ paper $\vd(t)$, $t=15$--$180$\,s.}
\label{tab:e3}
\begin{tabular}{@{}lccc@{}}
\toprule
\textbf{Metric} & \textbf{constrained (paper)} & \textbf{constrained (symm.)} & \textbf{conventional PD}\\
\midrule
crashed & no & no & no\\
time outside $\rho_q$ & 0\,\% & 0\,\% & 0.16\,\%\\
time outside $\rho_\omega$ & 6.0\,\% & \textbf{0\,\%} & 36.4\,\%\\
$\max|\varepsilon_1|$ / $\max|\varepsilon_2|$ & 0.93 / 1.21 & 0.72 / 0.77 & 1.02 / 1.25\\
tilt RMS & $0.43^\circ$ & $0.43^\circ$ & $0.52^\circ$\\
yaw RMS & $1.16^\circ$ & $1.06^\circ$ & $1.61^\circ$\\
max horizontal error & 12.3\,cm & 13.1\,cm & 14.9\,cm\\
\bottomrule
\end{tabular}
\end{table}

\begin{projbox}[conclusion of E3]
On the quadcopter the paper's qualitative claim holds again: with the same gains and limits, the constrained law cuts the
time outside the rate envelope from 36.4\,\% (PD) to 6.0\,\% (paper form) and to 0\,\% (symmetric form), and reduces the tilt
and yaw errors. The paper form's remaining violations are all at the lower yaw wall, the one-sided effect seen in the
spacecraft Monte-Carlo study B.
\end{projbox}

\begin{warn}[E3b: full paper disturbance]
At $1\times$ the paper's disturbance the yaw torque demanded exceeds what the X3's rotors can produce while hovering
(0.15 vs.\ 0.106\,N\,m): the mixer is saturated 97\,\% of the time and all three controllers are outside the envelopes
96--98\,\% of the time (none crashes; thrust has priority). This is the feasibility limit of \cref{sec:sens} in physical form:
no controller can satisfy a constraint the actuators cannot.
\end{warn}

% ---------------------------------------------------------------------
\subsection{E4: square trajectory}
% ---------------------------------------------------------------------

\resfig{x3_E4_square_trajectory.png}{\linewidth}{E4, $2\times2$\,m square at 2\,m altitude with minimum-jerk legs of 5\,s and 2\,s
holds. Left: 3-D path; middle: top view; right: position error (the first second is the take-off from the ground).}{fig:e4}

\begin{table}[H]
\centering\small
\caption{E4 metrics (after the take-off, $t\ge12$\,s).}
\label{tab:e4}
\begin{tabular}{@{}lcc@{}}
\toprule
\textbf{Metric} & \textbf{constrained (paper)} & \textbf{conventional PD}\\
\midrule
crashed & no & no\\
RMS / max position error & 4.03 / 7.90\,cm & 4.02 / 7.99\,cm\\
time outside $\rho_q$ & 0\,\% & 0\,\%\\
time outside $\rho_\omega$ & 4.7\,\% & 8.1\,\%\\
$\max|\varepsilon_1|$ & 0.30 & 0.36\\
\bottomrule
\end{tabular}
\end{table}

\begin{projbox}[conclusion of E4]
The paper's attitude law, with our outer loop, flies the square with 4\,cm RMS error and never leaves the attitude
envelope. Position tracking is set by the outer loop and is the same for both controllers. During the turns at the
corners the required roll and pitch rates exceed the hover-sized rate envelope ($\rho_{\omega T}=0.05$\,rad/s): the
constrained law halves the time outside it (4.7 vs.\ 8.1\,\%) but cannot remove it. For trajectory flight the rate envelope
should be widened to the rates the trajectory needs (feasibility again).
\end{projbox}

% ---------------------------------------------------------------------
\subsection{E5: control rate, delay and jitter}
% ---------------------------------------------------------------------

\begin{table}[H]
\centering\small
\caption{E5, disturbance $0.5\times$, 90\,s; time outside $\rho_q$ / $\rho_\omega$ in \%. No run crashed.}
\label{tab:e5}
\begin{tabular}{@{}lccc@{}}
\toprule
\textbf{Loop} & \textbf{constrained (paper)} & \textbf{constrained (symm.)} & \textbf{conventional PD}\\
\midrule
60\,Hz, 1 sample delay (nominal) & 0.0 / 6.6 & 0.0 / 0.0 & 0.4 / 36.2\\
100\,Hz, 1 sample delay & 0.0 / 6.8 & 0.0 / 0.0 & 0.4 / 36.7\\
250\,Hz, 1 sample delay & 0.0 / 7.1 & 0.0 / 0.0 & 1.5 / 36.9\\
60\,Hz, 2 samples delay & 18.0 / 90.6 & 20.7 / 90.6 & 0.4 / 36.3\\
60\,Hz, 2\,ms stamp jitter & 0.0 / 21.6 & 0.0 / 50.2 & 0.3 / 44.3\\
30\,Hz, 1 sample delay & 52.8 / 96.9 & 94.2 / 99.8 & 68.5 / 99.2\\
\bottomrule
\end{tabular}
\end{table}

\begin{warn}[What E5 tells us about deployment]
At 60\,Hz and faster the result is unchanged. Two things degrade the constrained laws far more than plain PD: a second
sample of delay (rate-envelope violations rise from 6.6 to 90.6\,\% for the paper form, while PD stays at 36\,\%) and
2\,ms of jitter on the time stamps (6.6 to 21.6\,\% for the paper form, 0 to 50.2\,\% for the symmetric form, 36 to 44\,\%
for PD). Both act through the body-rate estimate, which the barrier amplifies near the wall (even with the cap). At 30\,Hz no controller keeps the envelopes. In Gazebo the stamp must therefore be the simulation
time of the message (F12), and the pose rate should be at least 60\,Hz.
\end{warn}
